\chapter{Four saturation events}\label{ch:saturation}

\section{Two entries per event}

We propose that every actualization event makes two ledger entries \conjecture. The first
goes to the \emph{local} encoding surface, which is hot and finite: the transistor junction,
the qubit's substrate and barrier, the chromatin surface of a nucleus, or a forming
black-hole horizon. The second goes to the \emph{cosmic} horizon, which is cold and, on
observable timescales, never saturates. Information migrates from the first to the second at
a rate set by their temperature difference. A local surface can saturate; the cosmic one
cannot.

\section{The four thresholds}

\begin{center}\small
\begin{tabularx}{\textwidth}{@{}>{\raggedright\arraybackslash}p{0.14\textwidth}>{\raggedright\arraybackslash}X>{\raggedright\arraybackslash}X>{\raggedright\arraybackslash}p{0.27\textwidth}@{}}\toprule
domain & local surface & saturation observed as & status of the coefficient \\\midrule
gravity & horizon, area $A$ & collapse; $S=A/4\lP^2$ & \derived\ as a ratio, $2\pi/8\pi$ (Ch.~\ref{ch:blackholes}) \\
semiconductors & gate channel at $T_j$ & end of Dennard scaling, 2004--06 & set by $\kB T_j/q$ (Ch.~\ref{ch:cmos}); exponent $n$ \calibrated \\
quantum computing & junction barrier and substrate & substrate saturation point, $T_1^*=0.65\,T_{1,\mathrm{free}}$ & \calibrated\ (Ch.~\ref{ch:walls}) \\
cells & chromatin surface, at the cell's architecture position $P$ & to be measured (breach and the cancer region on the gauge) & floor $H(\varepsilon_0)$ from the measured holding energy (Part~\ref{part:4}) \measured \\\bottomrule
\end{tabularx}
\end{center}

\section{What is common, and what is not yet shown}

The common structure is real at the level of form. In each case:
\begin{enumerate}
\item there is a surface on which irreversible events are recorded;
\item the surface has a temperature: $\hbar\kappa/2\pi$, $T_j$, $\Delta/\kB$ or $T_{\mathrm{body}}$;
\item there is a floor below which a record cannot be written or kept, and a regime change
      when the system is driven against it.
\end{enumerate}
A stronger proposal is that each coefficient is a ratio of a topological
period to a geometric normalization, as $\tfrac14=2\pi/8\pi$ is for gravity \conjecture. For gravity the ratio is
derived (Chapter~\ref{ch:blackholes}). For the other three domains the
corresponding ratio has not been written down, and their coefficients are set from data. The four thresholds therefore stand as a
unified \emph{picture} with one derived case. Turning the
picture into a theory requires deriving at least one of the non-gravitational coefficients
($0.65$, a class exponent $n$, or a floor $\Hmin$) from surface geometry, before comparing it with
data \openprob.

\section{The bridge to the cell}

The biological application is the subject of Part~\ref{part:4}, and only its structure is needed here. The methylation state of a CpG is a binary record
written on chromatin, copied at every division. A cell's copy error $\varepsilon$, read on single molecules, is set against the thermal floor
$\varepsilon_0=1/(1+e^{E_{\rm hold}/\kB T})$ for a measured holding energy $E_{\rm hold}=3.41\,\kB T$, at the cell's architecture position $P$:
$A=H(\varepsilon)/(P\,H(\varepsilon_0))$ (IAM-A). On arrays, the same cell is read against the healthy cell of its own kind (Met-A). Both read
1 for a healthy cell, and both are per cell. Part~\ref{part:4} builds the instrument; the point here is only why a
physicist should expect a Landauer-type floor on a chromatin surface in the first place.
